\documentclass[a4paper]{article}

% 文章排版
\usepackage{indentfirst}
\setlength{\parindent}{0em}
\usepackage{autobreak}
\allowdisplaybreaks
\usepackage{parskip}
\usepackage{geometry}
\geometry{top=3cm,bottom=3cm,left=2cm,right=2cm}
\usepackage{anyfontsize}

% 数学物理宏包
\usepackage{amsmath}
\usepackage{physics}
\renewcommand\div\divisionsymbol
\DeclareMathOperator{\diag}{diag}
\DeclareMathOperator{\sgn}{sgn}
\newcommand{\trans}{\text{\textsf{T}}}
\usepackage{tabularray}
\UseTblrLibrary{amsmath}

\usepackage{xcolor}
\usepackage{fontspec}
\setmainfont{STIX Two Text}
\usepackage{unicode-math}
\unimathsetup{mathrm=sym,mathbf=sym,mathsf=sym,bold-style=ISO}
\setmathfont{STIX Two Math}

\usepackage[fontset=none]{ctex}
\setCJKmainfont{FZShuSong-Z01}[BoldFont=Source Han Sans CN,ItalicFont=FZKai-Z03]
\setCJKsansfont{Source Han Sans CN}

% 符号重定义
\AtBeginDocument{
    \let\uglyepsilon\epsilon
    \renewcommand\epsilon\varepsilon
    \let\uglyleq\leq
    \renewcommand\leq\leqslant
    \let\uglygeq\geq
    \renewcommand\geq\geqslant
}


% 题目
\title{\textbf{广义相对论（天文方向）\ \ \ \ 作业六}}
\author{国家天文台 张植竣}
\date{2025年10月21日}

\begin{document}

\maketitle

\begin{enumerate}
    \item 已知$F^{\mu\nu}$是一个反对称张量，$T_{\mu\nu}$是一个对称张量，证明$F^{\mu\nu}T_{\mu\nu}=0$。

    \textcolor{blue}{
    由于$F^{\mu\nu}$是反对称的，即$F^{\mu\nu} = -F^{\nu\mu}$，而$T_{\mu\nu}$是对称的，即$T_{\mu\nu} = T_{\nu\mu}$，我们有：
    \[
    F^{\mu\nu}T_{\mu\nu} = -F^{\nu\mu}T_{\nu\mu}
    \]
    另一方面，交换重命名$\mu$和$\nu$不应该改变求和的结果，因此我们也有：
    \[
    F^{\mu\nu}T_{\mu\nu} = F^{\nu\mu}T_{\nu\mu}
    \]
    这意味着：
    \[
    -F^{\nu\mu}T_{\nu\mu} = F^{\nu\mu}T_{\nu\mu}
    \]
    因此：
    \[
        F^{\mu\nu}T_{\mu\nu} = 0
    \]
    }

    \item 从Bianchi恒等式出发，证明Einstein张量的散度恒为零。

    \textcolor{blue}{
    Bianchi恒等式为：
    \[
        \nabla_\lambda R_{\mu\nu\rho\sigma} + \nabla_\rho R_{\mu\nu\sigma\lambda} + \nabla_\sigma R_{\mu\nu\lambda\rho} = 0
    \]
    Ricci张量定义为Riemann张量的如下收缩：
    \[
        R_{\nu\sigma} = R^{\mu}_{\ \nu\mu\sigma} = g^{\mu\rho} R_{\mu\nu\rho\sigma}
    \]
    Ricci标量定义为Ricci张量的如下收缩：
    \[
        R = R^{\mu}_{\ \mu} = g^{\mu\nu} R_{\mu\nu}
    \]
    Einstein张量的定义为：
    \[
        G_{\mu\nu} = R_{\mu\nu} - \frac{1}{2} R g_{\mu\nu}
    \]
    Einstein张量的散度为$\nabla^\mu G_{\mu\nu}$，我们需要证明$\nabla^\mu G_{\mu\nu} = 0$。 \\
    该证明需要对Bianchi恒等式做两次收缩：
    \[
        g^{\mu\rho} g^{\nu\sigma} \qty(\nabla_\lambda R_{\mu\nu\rho\sigma} + \nabla_\rho R_{\mu\nu\sigma\lambda} + \nabla_\sigma R_{\mu\nu\lambda\rho}) = 0
    \]
    对第一项：
    \[
        g^{\mu\rho} g^{\nu\sigma} \nabla_\lambda R_{\mu\nu\rho\sigma}
        = \nabla_\lambda (g^{\mu\rho} g^{\nu\sigma} R_{\mu\nu\rho\sigma})
        = \nabla_\lambda R
    \]
    对第二项：
    \begin{align*}
        g^{\mu\rho} g^{\nu\sigma} \nabla_\rho R_{\mu\nu\sigma\lambda}
        &= \qty(g^{\mu\rho} \nabla_\rho) \qty(g^{\nu\sigma} R_{\mu\nu\sigma\lambda}) \\
        &= \qty(g^{\mu\rho} \nabla_\rho) \qty(-g^{\nu\sigma} R_{\nu\mu\sigma\lambda}) \\
        &= -\nabla^\mu R_{\mu\lambda}
    \end{align*}
    对第三项：
    \begin{align*}
        g^{\mu\rho} g^{\nu\sigma} \nabla_\sigma R_{\mu\nu\lambda\rho}
        &= \qty(g^{\nu\sigma} \nabla_\sigma) \qty(g^{\mu\rho} R_{\mu\nu\lambda\rho}) \\
        &= \qty(g^{\nu\sigma} \nabla_\sigma) \qty(-g^{\mu\rho} R_{\mu\nu\rho\lambda}) \\
        &= -\nabla^\nu R_{\nu\lambda}
    \end{align*}
    将三项合并得：
    \[
        \nabla_\lambda R - \nabla^\mu R_{\mu\lambda} - \nabla^\nu R_{\nu\lambda} = 0
    \]
    即：
    \[
        \nabla^\mu R_{\mu\lambda} = \frac{1}{2} \nabla_\lambda R
    \]
    因此：
    \begin{align*}
        \nabla^\mu G_{\mu\nu}
        &= \nabla^\mu R_{\mu\nu} - \frac{1}{2} \nabla^\mu (g_{\mu\nu} R) \\
        &= \nabla^\mu R_{\mu\nu} - \frac{1}{2} g_{\mu\nu} \nabla^\mu R \\
        &= \nabla^\mu R_{\mu\nu} - \frac{1}{2} \nabla_\nu R \\
        &= 0
    \end{align*}
    这就证明了Einstein张量的散度恒为零。
    }

\end{enumerate}

\end{document}